\documentclass[11pt,a4paper]{article}

% ---------- Paquetes ----------
\usepackage[utf8]{inputenc}
\usepackage[T1]{fontenc}
\usepackage{lmodern}
\usepackage[spanish]{babel}
\usepackage{csquotes}
\usepackage[margin=2.5cm]{geometry}
\usepackage{graphicx}
\usepackage[backend=biber,style=ieee]{biblatex}
\usepackage[hidelinks]{hyperref}

\addbibresource{referencias.bib}

\begin{document}

% ---------- Carátula ----------
\begin{titlepage}
  \centering
  % \includegraphics[width=4cm]{capturas/logo-uteq.png}\\[1cm]
  {\large\bfseries UNIVERSIDAD TÉCNICA ESTATAL DE QUEVEDO\par}
  \vspace{0.3cm}
  {\large Facultad de Ciencias de la Computación\par}
  \vspace{0.2cm}
  {\large Carrera de Ingeniería de Software\par}
  \vspace{2cm}
  {\Large\bfseries FORM-IC-01\par}
  \vspace{0.4cm}
  {\Large Investigación individual: arquitecturas distribuidas en la industria\par}
  \vspace{0.6cm}
  {\large\itshape Caso de estudio: la arquitectura de microservicios orientada a dominios (DOMA) de Uber\par}
  \vfill
  \begin{tabular}{rl}
    \textbf{Asignatura:} & Aplicaciones Distribuidas \\
    \textbf{Estudiante:} & Mera Jumbo Allison Michelle \\
    \textbf{Docente:}    & Dr. Gleiston Cicerón Guerrero Ulloa \\
    \textbf{Nivel:}      & Séptimo nivel, paralelo A \\
  \end{tabular}
  \vspace{1.5cm}

  {\large Quevedo, noviembre de 2026\par}
\end{titlepage}

% ---------- Análisis (máximo 1 página) ----------
\section*{1. Descripción del caso}
% Qué empresa, qué sistema y a qué escala, con cifras citadas.
% Ejemplo de cita: \cite{uber2020doma}

\section*{2. Por qué eligió este estilo y qué tenía antes}
% El monolito inicial, los problemas al crecer y el paso a microservicios y luego a DOMA.

\section*{3. Beneficios obtenidos}
% Datos medibles que da la fuente.

\section*{4. Desafíos y cómo los resolvió}
% Complejidad de miles de microservicios, depuración, etc., y cómo DOMA los atacó.

\section*{5. Aplicación al PFC}
% Qué se tomaría para el sistema del CRTV y qué no, y por qué.

% ---------- Referencias ----------
\newpage
\printbibliography[title={Referencias}]

% ---------- Anexo: evidencias de autoría ----------
\newpage
\section*{Anexo A. Evidencias de autoría}

\subsection*{A.1 Captura de la fuente principal anotada}
% Cuando tengas la imagen en capturas/, quita los % de las líneas de abajo:
% \begin{figure}[h]
%   \centering
%   \includegraphics[width=\textwidth]{capturas/uber-blog-anotado.png}
%   \caption{Fragmentos subrayados del artículo oficial de Uber.}
% \end{figure}

\subsection*{A.2 Ejemplo de primera mano}
% Un sistema que uses o conozcas en el que reconozcas microservicios.

\subsection*{A.3 Declaración de uso de inteligencia artificial}
% Qué herramienta usaste, para qué y qué partes son de tu autoría.

\end{document}